\section{Protocolo de avaliação}
\label{sec:protocolo}

\subsection{Objetivo e desenho comparativo}

A avaliação responderá se a abordagem desenvolvida oferece uma relação entre custo e qualidade mais favorável do que a execução de um modelo em seu harness padrão, sob uma restrição de latência. A referência prática será uma configuração documentada do modelo forte no harness selecionado, como Claude Code ou Codex, preservando os recursos nativos de execução. Esses nomes identificam candidatos a referência; o protocolo não pressupõe equivalência entre eles nem exige usar um modelo em um produto que não o suporte. Cada conclusão será vinculada à combinação efetivamente avaliada de modelo, harness e parâmetros.

O contraste principal será entre sistemas completos. Cada condição receberá as mesmas tarefas e estados iniciais, com permissões, acesso a ferramentas externas, testes disponíveis e tetos de recursos comparáveis. A organização interna do trabalho poderá variar, pois constitui parte da intervenção. Comparações que isolem um componente manterão os demais fatores fixos; quando isso não for possível, a diferença será atribuída ao conjunto da configuração. A ordem das execuções será aleatorizada em blocos por tarefa para reduzir a associação entre condição e horário de acesso ao serviço.

\subsection{Condições e contrastes}

Serão avaliadas cinco condições:
\begin{enumerate}
    \item \textbf{Modelo forte no harness padrão:} referência principal de resolução, custo e latência, com versão e configuração de execução registradas.
    \item \textbf{Modelo barato no harness de referência compatível:} controle para medir o benefício e a perda de resolução obtidos pela substituição do modelo. O harness será mantido quando houver suporte; diferenças necessárias de harness serão documentadas e impedirão atribuir o contraste exclusivamente ao modelo.
    \item \textbf{Modelo forte com economia de tokens:} o mesmo modelo e harness da primeira condição, acrescidos de técnicas explicitamente selecionadas para reduzir consumo. As técnicas serão definidas na investigação dos mecanismos e congeladas antes da avaliação final. Seu efeito será medido em custo monetário, resolução e tempo; reduzir tokens, isoladamente, não caracterizará benefício.
    \item \textbf{Modelo barato com escalonamento simples:} execução inicial pelo modelo barato e intervenção do forte segundo uma regra fixa, baseada apenas em sinais disponíveis durante a execução. A regra, os limites e a transferência de contexto serão documentados.
    \item \textbf{Abordagem desenvolvida:} configuração selecionada na etapa de desenvolvimento, que poderá empregar um ou mais modelos e mecanismos de economia, verificação ou coordenação.
\end{enumerate}

A quinta condição será comparada à primeira para responder à pergunta principal. Os contrastes com a segunda, a terceira e a quarta avaliarão se o benefício também pode ser obtido por alternativas mais simples. Superar apenas a referência padrão não demonstrará superioridade sobre esses controles. O contraste entre a primeira e a terceira condições medirá especificamente a contribuição das técnicas de economia aplicadas ao modelo forte.

As otimizações nativas da referência serão preservadas; a terceira condição acrescentará intervenções identificáveis sobre essa base. Todas as condições terão acesso à mesma informação externa permitida, sem testes reservados ou soluções de referência. Configurações ajustáveis receberão oportunidades documentadas de ajuste nas tarefas de desenvolvimento. Ablações corresponderão aos componentes efetivamente adotados e às alegações sobre seus efeitos. Um controle adicional de repetição simples será incluído se repetição for um mecanismo central, distinguindo o efeito da política do benefício de gastar mais inferência \cite{Kapoor2024AIAgentsMatter}.

\subsection{Contrastes de contexto, composição e processo}

O desenvolvimento examinará conjuntamente modelo e estratégia de contexto em um conjunto limitado de variantes. Para a composição por atividade, o piloto cruzará, quando tecnicamente compatíveis, dois modelos nas etapas de seleção de contexto e resolução. As combinações com o mesmo modelo nas duas etapas serão controles para as combinações heterogêneas, mantendo representações, ferramentas, amostragem de candidatos e verificação comparáveis. Uma configuração de execução direta permitirá verificar se o benefício decorre da separação em etapas. Para isolar o resolvedor, serão reutilizados os mesmos contextos selecionados; para isolar o seletor, será mantido o resolvedor. O custo da seleção integrará a configuração completa, mesmo quando seus artefatos forem reutilizados nos contrastes diagnósticos.

Uma composição fixa não será denominada roteamento dinâmico. Se houver escolha condicional de modelos, a política será comparada com atribuições fixas e com o escalonamento simples. Analogamente, uma política de ativação de exploração dedicada será comparada com as alternativas de sempre executar e sempre dispensar essa etapa, mantendo as possibilidades de inspeção durante o reparo. Gatilhos dependerão apenas de informações acessíveis ao agente, como descrição da issue, evidências recuperadas e falhas observadas. Soluções de referência, arquivos do patch oficial e resultados dos testes reservados não poderão orientar decisões operacionais.

Intervenções de contexto serão descritas por sua fonte, granularidade, extensão, ordenação, atualização e regras de compartilhamento. Se forem adotadas skills, serão registrados conteúdo, versão, mecanismo de carregamento e critério de seleção; os controles distinguirão ausência de skill, disponibilização fixa e seleção condicional, conforme a alegação investigada. Para isolamento de validação, o contraste especificará quais informações são fornecidas ao revisor e ao editor, mantendo evidências executáveis e orçamento comparáveis. Tais contrastes serão exigidos apenas para mecanismos efetivamente adotados, evitando um produto cartesiano de todas as variantes. O piloto selecionará as comparações confirmatórias antes da avaliação reservada.

\subsection{Benchmarks}

O SWE-bench consolidou a avaliação de agentes que recebem uma issue e um repositório e devem produzir um patch verificável por execução de testes \cite{Jimenez2024SWEbench}. O SWE-bench-Live será o conjunto principal para esse tipo de tarefa: sua coleção atualizável reúne issues mais recentes, de um número maior de repositórios, com ambientes executáveis \cite{Zhang2025SWEbenchLive}. A atualidade reduz uma fonte potencial de exposição prévia das tarefas, mas não comprova ausência de contaminação. Tampouco uma diferença de pontuação entre conjuntos demonstra, isoladamente, memorização: a dificuldade e a composição das tarefas também podem variar.

O Multi-SWE-bench complementará a avaliação por incluir 1.632 tarefas verificadas em Java, TypeScript, JavaScript, Go, Rust, C e C++ \cite{MultiSWEbench2025}. O estudo desse conjunto relata diferenças de desempenho entre linguagens e entre métodos de agente. O SWE-PolyBench também documenta desempenho desigual e distribuição não uniforme de tarefas por linguagem, reforçando a necessidade de apresentar resultados desagregados \cite{Rashid2025SWEPolyBench}. Python aparece nas comparações publicadas do Multi-SWE-bench por meio de tarefas de outro conjunto; não será contado como uma das sete linguagens que compõem suas 1.632 instâncias.

O Terminal-Bench será a terceira avaliação, voltada à execução de tarefas por agentes em ambientes de linha de comando. Sua publicação acadêmica apresenta o Terminal-Bench 2.0 e um protocolo de verificação por tarefa \cite{Merrill2026TerminalBench}. O projeto mantém versões posteriores; a versão 4.0 alterou tarefas, verificadores e recursos de execução \cite{Marten2026TerminalBench4}. A versão e o identificador do conjunto serão fixados antes dos experimentos. Como suas tarefas abrangem também domínios além da engenharia de software, seus resultados serão interpretados como evidência complementar sobre atuação em terminal, e não como taxa de resolução de issues.

A inclusão de cada instância exigirá ambiente reproduzível, instruções suficientes e um verificador aplicável. Auditorias do SWE-bench Verified e do SWE-bench Pro encontraram, respectivamente, problemas de exposição de soluções e de adequação dos testes, e problemas de qualidade em tarefas do conjunto Pro \cite{OpenAI2026SWEbenchVerified,OpenAI2026SignalNoise}. Esses achados motivam a auditoria e o registro de problemas nas instâncias usadas; não serão extrapolados automaticamente para o SWE-bench-Live, o Multi-SWE-bench ou o Terminal-Bench. Não haverá agregação das pontuações dos três benchmarks em um escore único.

\subsection{Execução, entrega e registro de falhas}

Cada execução começará em ambiente isolado no estado especificado pela tarefa. Nos benchmarks de issues, o agente editará os arquivos do repositório; após o encerramento, um coletor comum produzirá o patch em relação ao commit inicial, removendo apenas artefatos previamente especificados. Serão preservados a cópia de trabalho final e o patch submetido. O coletor, suas regras e a versão do avaliador serão comuns às condições, evitando introduzir exigências diferentes de serialização da solução \cite{Zheng2026ClawSWEBench}. No Terminal-Bench será aplicado o contrato de avaliação da versão selecionada.

O encerramento ocorrerá por entrega, abandono ou esgotamento de um limite de tempo, gasto ou passos. Tentativas internas e intervenções auxiliares integrarão a mesma execução e consumirão seu orçamento. Repetições experimentais começarão novamente no estado inicial e serão registradas separadamente; o melhor resultado entre repetições não será usado como resultado de uma execução comum. O avaliador reservado será acionado após o encerramento e não fornecerá sinais para novas tentativas do agente.

Nas condições que empregarem verificação interna, serão especificados o verificador, as evidências acessíveis e a regra que transforma seu resultado em aceitação, reparo, nova tentativa, escalonamento ou encerramento. Os registros distinguirão intervenções que apenas julgam uma solução daquelas que orientam ou executam alterações. Limiares e regras serão definidos no desenvolvimento e congelados antes da avaliação final. Chamadas de revisão e testes internos integrarão custo e latência, permitindo avaliar a política completa, inclusive quando detectar uma falha não resultar em reparo bem-sucedido.

Ausência de entrega, patch inválido, falha nos testes e esgotamento de recursos contarão como insucesso quando decorrerem da configuração avaliada. A causa e a etapa serão registradas: localização, edição, entrega, teste, verificação, escalonamento ou parada. Falhas comprovadamente externas, como indisponibilidade de serviço ou ambiente corrompido, terão tratamento separado, segundo critérios e limite de reexecução fixados no piloto e aplicados igualmente às condições. Os registros originais e os gastos dessas ocorrências serão preservados, inclusive quando não entrarem na estimativa principal. Nenhuma falha será descartada apenas por produzir resultado desfavorável.

Serão registrados versões de modelos, harnesses e adaptadores, parâmetros de geração, prompts, permissões, limites, datas, chamadas, uso de tokens, ações de ferramenta, testes internos, decisões de escalonamento e motivos de encerramento. Procedência e datas das tarefas serão preservadas. A separação entre desenvolvimento e avaliação final impedirá que os resultados reservados sejam usados para selecionar configurações.

Quando houver recuperação estrutural ou índices, sua construção usará exclusivamente o estado inicial permitido do repositório. Atualizações acompanharão as alterações realizadas durante a execução, sem acesso a commits de solução ou testes reservados. Serão registrados versão do índice, consultas, trechos retornados e conteúdo efetivamente fornecido a cada etapa. Para distinguir aquecimento de economia recorrente, custos e tempos de preparação serão apresentados separadamente e incorporados ao cenário de primeira utilização. Qualquer amortização entre tarefas explicitará o número de usos, a validade do índice e os custos de atualização; o compartilhamento de cache ou índice entre condições não poderá favorecer apenas uma delas.

\subsection{Métricas}

Para um benchmark $b$, seja $I_b$ o conjunto de tarefas, com tamanho $n_b$, e $K$ o número comum de repetições experimentais por tarefa e condição. Em cada execução $r$, $s_{i,c,r}=1$ indicará aprovação pelo verificador oficial e $s_{i,c,r}=0$ indicará insucesso. A taxa de resolução será
\begin{equation}
\widehat{R}_{b,c}=\frac{1}{n_bK}\sum_{i\in I_b}\sum_{r=1}^{K}s_{i,c,r}.
\end{equation}
Essa taxa mede correção segundo o protocolo do benchmark. Nos conjuntos de issues, serão exigidos os testes de correção e de ausência de regressão previstos pelo respectivo avaliador \cite{Jimenez2024SWEbench}. A taxa de aplicação de patches e a fração de testes aprovados serão diagnósticos, sem converter aprovação parcial em resolução completa \cite{Zhou2026FeatureBench}.

O custo de API de uma execução incluirá planejamento, execução, revisão, resumos, escalonamentos e tentativas internas, inclusive chamadas malsucedidas cobradas. Para cada chamada $j$, sejam $u_{j,k}$ os tokens da categoria de cobrança $k$, $p_{m_j,k}$ o preço por milhão de tokens do modelo empregado e $f_j$ eventuais tarifas fixas:
\begin{equation}
C^{\mathrm{API}}_{i,c,r}=\sum_{j\in J_{i,c,r}}\left(\sum_k\frac{u_{j,k}p_{m_j,k}}{10^6}+f_j\right).
\end{equation}
As categorias de entrada, cache e saída serão mutuamente exclusivas. Moeda, data, fonte das tarifas e chamadas sem informação de cobrança serão registradas; dados ausentes não serão tratados como custo zero. O custo calculado será confrontado com a cobrança quando disponível. Para referências acessadas por assinatura, eventual valoração por tarifas de API será identificada como estimativa, separada do desembolso da assinatura, e só será utilizada se houver dados suficientes de uso e preços verificáveis.

A medida monetária principal será o custo médio por tarefa, incluindo falhas, $\overline{C}_{b,c}=\sum_{i,r}C^{\mathrm{API}}_{i,c,r}/(n_bK)$. A medida complementar será o custo por resolução, $\sum_{i,r}C^{\mathrm{API}}_{i,c,r}/\sum_{i,r}s_{i,c,r}$, indefinido quando não houver resolução. O numerador incluirá todas as execuções. Infraestrutura e ferramentas serão informadas separadamente; uma redução apenas no custo de API não será apresentada como redução comprovada do custo operacional total. Tokens, cache, número de chamadas e participação do modelo forte serão diagnósticos de consumo.

A latência será $T_{i,c,r}=t^{\mathrm{fim}}_{i,c,r}-t^{\mathrm{inicio}}_{i,c,r}$, do fornecimento da tarefa ao agente até a entrega ou o encerramento, incluindo esperas, ferramentas, testes internos e coordenação. O tempo de avaliação externa será registrado à parte. A medida principal será a média $\overline{T}_{b,c}$ sobre todas as execuções, incluindo falhas e esgotamentos de tempo; também serão apresentados mediana, percentil 95, distribuição e frequência de encerramentos por limite. A duração de trabalhos paralelos será medida pelo tempo decorrido, e seu custo incluirá todas as chamadas.

A confiabilidade da conclusão interna será descrita por
\begin{equation}
\widehat{F}_{b,c}=\frac{\sum_{i,r}d_{i,c,r}(1-s_{i,c,r})}{\sum_{i,r}d_{i,c,r}},
\end{equation}
em que $d_{i,c,r}=1$ indica declaração de tarefa concluída. O indicador ficará indefinido na ausência dessas declarações. Abandonos e encerramentos sem solução serão apresentados separadamente.

Esse indicador mede a fração de declarações de conclusão rejeitadas pelo avaliador oficial. Seu denominador difere do de uma taxa de falsos positivos condicionada a candidatos incorretos; ele não estima o erro de cada julgamento intermediário nem a correção semântica além do benchmark. Sua comparação entre condições será acompanhada da frequência de declarações de conclusão e da taxa de resolução, para que declarar menos conclusões não seja interpretado isoladamente como melhoria do sistema.

As expressões acima descrevem as tarefas avaliadas com pesos iguais. Uma estimativa para uma população com probabilidades desiguais de inclusão empregará os pesos do plano amostral, também nas estimativas de custo e tempo. Repetições da mesma tarefa não serão tratadas como novas tarefas independentes.

Os diagnósticos incluirão custo, tempo e chamadas por atividade, modelo empregado, extensão do contexto, etapas ativadas e descartadas e número de candidatos gerados e julgados. Esses registros permitirão distinguir economia local de redução do gasto total: contexto menor pode gerar mais consultas, e uma etapa adicional pode evitar reparos posteriores. Métricas de localização baseadas nos arquivos da solução de referência serão calculadas somente após a execução e interpretadas como aproximações, pois soluções alternativas podem modificar outros arquivos. Elas não substituirão a resolução final nem definirão retrospectivamente o contexto ideal disponível ao agente.

\subsection{Critérios de benefício e aceitação}

Qualidade, custo e tempo serão analisados conjuntamente, sem redução a um escore único. Para uma configuração $c$ e a referência forte $c_0$, no mesmo benchmark, definem-se a diferença de resolução $\Delta R_c=R_{b,c}-R_{b,c_0}$, a economia relativa $E_c=1-\overline{C}_{b,c}/\overline{C}_{b,c_0}$ e a razão de latência $L_c=\overline{T}_{b,c}/\overline{T}_{b,c_0}$, para denominadores positivos. As estimativas desses contrastes respeitarão o pareamento por tarefa.

O protocolo adota dois cenários suficientes de benefício econômico: redução de custo com taxa de resolução preservada ou superior ($E_c>0$ e $\Delta R_c\geq0$); e economia de pelo menos 90\% com perda de até cinco pontos percentuais ($E_c\geq0{,}90$ e $\Delta R_c\geq-0{,}05$). Em ambos, o limite operacional inicial será $L_c\leq2$. O segundo cenário corresponde, por exemplo, a reduzir o custo médio de US\$~5 para US\$~0,50 por tarefa, admitindo uma queda de resolução de 60\% para 55\%. Esses valores expressam critérios de projeto, não resultados experimentais nem limites estabelecidos pela literatura.

Os cenários não impõem economia mínima de 90\% a toda configuração nem autorizam perda de cinco pontos para qualquer economia. Outras combinações serão apresentadas como trocas entre dimensões, sem aceitação automática ou interpolação de uma tolerância não especificada. Configurações dominadas, isto é, sem vantagem em nenhuma das três dimensões e piores em pelo menos uma, serão identificadas; configurações não dominadas não serão, por esse fato, consideradas aceitáveis. O limite de latência se aplica à média e será acompanhado pelo percentil 95 e pelo teto individual de execução, evitando interpretar aprovação média como garantia para cada tarefa.

Os requisitos de cada cenário deverão ser satisfeitos pela mesma configuração na mesma comparação. A análise empregará intervalos de confiança com cobertura conjunta de 95\% para os contrastes confirmatórios, com procedimento de construção declarado antes da avaliação final. A aceitação exigirá que os limites inferiores sustentem os requisitos de resolução e economia e que o limite superior da razão de latência não exceda dois. Uma estimativa pontual favorável será considerada inconclusiva quanto à aceitação se sua incerteza não permitir sustentar esses requisitos. A ausência de diferença significativa não demonstrará preservação de qualidade. O cenário de perda limitada permitirá analisar não inferioridade com margem de cinco pontos, mas a aceitação exigirá também satisfazer o requisito de economia correspondente e a restrição de tempo.

Os contrastes com os demais controles serão apresentados nas mesmas dimensões. Um benefício frente ao harness padrão estabelecerá vantagem apenas frente àquela referência; vantagens adicionais atribuídas à complexidade da abordagem dependerão das comparações com economia de tokens, substituição do modelo e escalonamento simples.

\subsection{Piloto, dimensionamento e análise final}

O piloto será realizado em tarefas de desenvolvimento e terá cinco entregas: estimativas de custo e duração por condição; variabilidade entre repetições e frequência de resultados discordantes entre condições na mesma tarefa; validação da entrega e dos registros; diagnóstico de falhas externas e sinais observáveis de dificuldade; e cenários de orçamento para a avaliação final. Esses dados orientarão o tamanho da amostra, o número de repetições, os limites de gasto, tempo e passos e o tratamento de falhas externas. A regra do escalonamento simples será ajustada apenas nessa etapa. Tentativas internas, repetições experimentais e passos terão limites distintos.

A triagem de variantes poderá usar tarefas selecionadas por dificuldade histórica, conforme a distinção estabelecida na introdução \cite{Ndzomga2026EfficientBenchmarking}. Seus resultados orientarão o desenvolvimento, sem serem tratados como estimativas do benchmark inteiro. A avaliação confirmatória empregará tarefas reservadas e um plano de seleção previamente registrado. Versões, recortes e alocação entre benchmarks permanecerão decisões de dimensionamento; a seleção atual dos três conjuntos será mantida até essa definição.

O dimensionamento considerará a precisão dos contrastes pareados de resolução, custo e tempo, a dependência entre tarefas de um mesmo repositório e a variabilidade entre repetições. Antes da avaliação final serão congelados as configurações, os controles, as ablações pertinentes, os identificadores das tarefas, o plano amostral, os tetos de recursos, o método de construção dos intervalos e a regra de decisão conjunta. O piloto informará a viabilidade estatística e financeira; os critérios de aceitação não serão relaxados apenas para acomodar resultados favoráveis a uma configuração. Se houver precisão insuficiente, serão reduzidas variantes exploratórias, ampliada a avaliação ou limitada a conclusão, sem declarar benefício demonstrado com base apenas na estimativa pontual.

Os resultados serão apresentados separadamente por benchmark. No Multi-SWE-bench serão também desagregados por linguagem e repositório, acompanhados do tamanho dos grupos; esses contrastes serão descritivos quando a amostra não sustentar inferência própria. Patches, logs e diagnósticos permitirão examinar onde custo, tempo e resolução se alteraram. Uma auditoria complementar de requisitos não cobertos pelos testes oficiais só sustentará alegações adicionais se seu procedimento, amostra e critérios forem definidos previamente, com revisão sem identificação da configuração. A avaliação principal permanecerá vinculada ao verificador oficial.
